Fix a coordinate \(k=(a,y)\) and let \(C_k\) denote the corresponding cell count under a multinomial law with parameter \(q\).  The elementary first and second factorial-moment identities are obtained by expanding the count into indicators:
\[
 \sum_{c\in\mathcal I_m}\frac{c_k}{m}B_c^{(m)}(q)=q_k,
 \qquad
 \sum_{c\in\mathcal I_m}\frac{c_k(c_k-1)}{m(m-1)}B_c^{(m)}(q)=q_k^2,
\]
where the second identity uses \(m\geq2\).

Let \(b=g_m(q^*)\) and \(\nu\in\mathcal F_{m,\epsilon}(b)\).  The fiber equations say \(\int B_c^{(m)}(q)\,d\nu(q)=B_c^{(m)}(q^*)\) for every \(c\).  Applying the two displayed linear combinations gives
\[
 \int q_k\,d\nu(q)=q_k^*,
 \qquad
 \int q_k^2\,d\nu(q)=(q_k^*)^2.
\]
Hence \(\int(q_k-q_k^*)^2\,d\nu=0\), so \(q_k=q_k^*\) almost surely.  Doing this for all four coordinates yields \(q=q^*\) almost surely and therefore \(\nu=\delta_{q^*}\).  The fiber is the claimed singleton, and its causal image is \(L_m(b)=U_m(b)=f(q^*)\).